\begin{explanationbox}
	\textbf{\textcolor{CExplanation}{一句话直觉}}

	函数图像相对于坐标轴或原点的对称，本质是点的坐标符号取反或互换，从而得到新函数表达式。

	\medskip
	\textbf{\textcolor{CExplanation}{核心拆解}}
	\begin{description}[style=nextline, leftmargin=2.8em, labelsep=0.5em, itemsep=0.45em]
		\item[\textbf{要点一（坐标对应变换）：}] 图像上任意点 $(x, y)$ 经过对称后得到新点 $(x', y')$，新点的坐标满足特定的转换规则。例如关于 $x$ 轴对称时 $(x',y')=(x,-y)$。
		\item[\textbf{要点二（解析式替换）：}] 要求对称后的函数解析式，只需将新点坐标代入原函数关系，解出 $y'$ 关于 $x'$ 的表达式，再换回常用变量 $x,y$。
	\end{description}

	\medskip
	\textbf{\textcolor{CExplanation}{几何本质（点坐标变换）}}

	对称变换是平面上的等距变换，保持图形全等。图像对称后，点与对称轴/中心的距离不变，仅位置镜像。

	\medskip
	\textbf{\textcolor{CExplanation}{代数意义（变量符号替换）}}

	每个对称变换都对应一个简单的代数替换：关于 $x$ 轴用 $y$ 替换为 $-y$；关于 $y$ 轴用 $x$ 替换为 $-x$；关于原点同时替换 $x$ 和 $y$；关于 $y=x$ 交换 $x$ 与 $y$。

	\medskip
	\textbf{\textcolor{CExplanation}{考点价值（图像变换应用）}}
	\begin{itemize}[leftmargin=*, itemsep=0.5em, topsep=0.4em]
		\item \textbf{考法一（求对称解析式）：} 已知原函数，利用对称规则直接写出新函数解析式。
		\item \textbf{考法二（判断对称性质）：} 检验函数是否满足 $f(x)=f(-x)$（偶函数）或 $f(x)=-f(-x)$（奇函数）等。
	\end{itemize}

	\medskip
	\textbf{\textcolor{CExplanation}{顿悟点}}

	对称变换 = 点的坐标替换。记住每种对称对应的坐标变换规则，就能直接写出函数表达式。

	\medskip
	\textbf{\textcolor{CExplanation}{使用场景}}
	\begin{itemize}[leftmargin=*, itemsep=0.5em, topsep=0.4em]
		\item \textbf{场景一（函数图像绘制）：} 利用对称变换从已知函数图像快速得到相关图像。
		\item \textbf{场景二（反函数求解）：} 通过求原函数关于 $y=x$ 的对称图像得到反函数。
	\end{itemize}
\end{explanationbox}
